%!TEX root = ../main.tex
% =============================================================================
%  四、问题 1：确定性条件下的计划购电模型      目标 4 页
%  评审在问：确定性版本做对了吗？
%  本节含：① 电量口径与 LP 全式（目标／平衡含弃光松弛／SOC 递推／SOC 界／功率界／日周期）
%          ② 性质引理（充放互斥、能量恒等式）与证明
%          ③ 对偶分析：水价值 ν_t、并网点边际价值 α_t、阈值律与死区宽度
%          ④ 无储能基线 C-01、成本分解恒等式、容量/功率扫描、效率口径极差
%  图名：fig03_q1_dual_price.pdf / fig10_strategy_mechanism.pdf
%  单位纪律：功率 kW、电量 kWh、电价 元/kWh、费用 元；Δt=1/6 h；P̄Δt=833.3333 kWh
% =============================================================================

\section{问题 1：确定性条件下的计划购电模型}\label{chap:q1}

\subsection{模型}

\textbf{口径。} 采用电量口径：决策变量均以电量（kWh）为单位，电价 \(p_t\)（元/kWh）与电量相乘即
费用，\textbf{不再乘 \(\Delta t\)}；负载与光伏（kW）进模型前乘以 \(\Delta t\) 化为电量。
电价与负载逐日相同、光伏给出一日预测，故只需求解一个典型日（附件 1），信息结构为“当日真值已知”。
储能参数：额定容量 \(12\,000\)~kWh，运行区间 \([\underline S,\overline S]=[1200,10\,800]\)~kWh
（可用摆幅 9600~kWh），额定功率 \(\overline P=5000\)~kW，单边效率 \(\eta_c=\eta_d=0.9\)（假设 A2）。

决策变量为 \((E_t,c_t,d_t,s_t,S_t)\)，\(t\in\mathcal T\)，\(|\mathcal T|=144\)。模型为

\begin{subequations}\label{eq:q1-model}
\begin{align}
  \min_{E,c,d,s,S}\quad & C^{\text{plan}}=\sum_{t}p_tE_t \label{eq:q1-obj}\\
  \text{s.t.}\quad
  & E_t+(PV_t-L_t)\Delta t+d_t-c_t-s_t=0, && \forall t \label{eq:q1-balance}\\
  & S_t=S_{t-1}+\eta_c c_t-d_t/\eta_d, && \forall t \label{eq:q1-soc}\\
  & \underline S\le S_t\le\overline S,\quad 0\le c_t,d_t\le \overline P\Delta t, && \forall t \label{eq:q1-socbound}\\
  & E_t\ge0,\quad s_t\ge0, && \forall t \label{eq:q1-nonneg}\\
  & S_0=S_{144}=S^{init}=6000\ \text{kWh}. \label{eq:q1-periodic}
\end{align}
\end{subequations}

式~\eqref{eq:q1-balance} 即“微网供给不低于小区负载”的代数表达，\((PV_t-L_t)\Delta t\) 为净负荷电量
（盈余为负）。光伏盈余可能超过充电功率与负载之和，故以\textbf{等式加松弛 \(s_t\ge0\)} 写出以保证
可行；典型日的日内盈余低于摆幅 9600~kWh 与单槽吸纳能力 \(\overline P\Delta t=833.3333\)~kWh，
故最优 \(s_t\equiv0\)。同一典型日上“含松弛等式／不含松弛等式／不等式”三种写法的
最优值完全相同（均为 \(35\,126.9486\)~元），保留 \(s_t\) 的意义在于实际日期序列
（窗口内光伏峰值 \(10\,216.2000\)~kW）的可行性。式~\eqref{eq:q1-periodic} 即题面“两端储电量相同”，
本文进一步取两者均等于附件给定的初值 6000~kWh。

模型规模为 \(4\times144=576\) 个变量、288 条约束（显式保留 \(s_t\) 时为 720 个变量），是连续 LP；
充放互斥由模型结构自动成立，\textbf{无需 0--1 变量}。

\subsection{两条结构性质}

\begin{proposition}[充放互斥]\label{prop:no-simul}
若 \(\eta_c\eta_d<1\)，则最优解中不存在 \(c_t,d_t\) 同时为正的时段。
\end{proposition}

\begin{proof}
设某 \(t\) 有 \(c_t>0,d_t>0\)。取 \(\delta=\min\{c_t,\ d_t/(\eta_c\eta_d)\}>0\)，作
\(c_t'=c_t-\delta\)、\(d_t'=d_t-\eta_c\eta_d\delta\ge0\)，其余不变。由式~\eqref{eq:q1-soc}，
储电量变化 \(\eta_c(-\delta)+\eta_c\eta_d\delta/\eta_d=0\)，故储电量轨迹与全部界约束不变，新解可行；
由式~\eqref{eq:q1-balance}，购电量变化
\(\Delta E_t=-\delta+\eta_c\eta_d\delta=-\delta(1-\eta_c\eta_d)<0\)，目标下降
\(p_t\delta(1-\eta_c\eta_d)>0\)，与原解最优矛盾。
\end{proof}

命题~\ref{prop:no-simul} 只用 \(\eta_c\eta_d<1\)、\(p_t>0\) 与平衡式结构，对问题 1--4 的全部变体
一致成立。数值验证：同时充放时段数 0，储电量递推残差 \(9.09\times10^{-13}\)~kWh，
充放比偏差 \(4.44\times10^{-16}\)（审计 A5）。

\begin{lemma}[能量恒等式]\label{lem:energy}
逐槽由式~\eqref{eq:q1-balance} 解出 \(E_t\) 并对 \(t\) 求和，再用式~\eqref{eq:q1-soc} 消去充放电量，得
\begin{equation}
  \sum_{t}E_t=\sum_{t}(L_t-PV_t)\Delta t+\sum_{t}c_t-\sum_{t}d_t
  \ \xrightarrow{\ S_{144}=S_0\ }\
  \sum_{t}(L_t-PV_t)\Delta t+\bigl(1-\eta_c\eta_d\bigr)\sum_{t}c_t .
  \label{eq:q1-energy}
\end{equation}
\end{lemma}

引理~\ref{lem:energy} 有三点用途。① \textbf{它是“恒价无套利”的代数根源}：若日内电价恒为
\(\bar p\)，则目标 \(\bar p\sum_tE_t\) 只随循环量 \(\sum_tc_t\) 变化，与充放时序无关，储能\emph{没有}
时段套利价值；反之，题面背景段所指的“低充高放”只有在日内电价随时刻变化时才有对象，这从代数上支持
假设 A11（附件 1 的 144 点曲线逐日重复）。② \textbf{它是零成本的可行性回归校验}：逐槽核验
式~\eqref{eq:q1-energy}，典型日两侧之差在 \(10^{-4}\)~kWh 量级（即四位小数取整量级）。
③ \textbf{它分离了两问的价值来源}：问题 1 的节省来自改变 \(\sum_tE_t\) 的\emph{时段结构}（套利），
问题 2 的安全边际来自改变其\emph{分布}（降低高价缺口的期望），两者可叠加但不可互替。

\subsection{对偶分析与阈值型最优策略}

记 \(\nu_t\) 为式~\eqref{eq:q1-soc} 的影子价格，即\textbf{单位储电量的边际价值}（水价值，元/kWh）；
\(\alpha_t\) 为式~\eqref{eq:q1-balance} 的影子价格（并网点电量的边际价值，元/kWh，取非负约定）。
在储电量界与功率界均不活跃时，对 \(c_t,d_t\) 的平稳性条件给出互补松弛对

\begin{equation}
  z^{c}_t=\alpha_t-\eta_c\nu_t\ge0,\quad c_tz^{c}_t=0;\qquad
  z^{d}_t=-\alpha_t+\nu_t/\eta_d\ge0,\quad d_tz^{d}_t=0 .
  \label{eq:q1-kkt}
\end{equation}

其经济含义为：多充 1~kWh 的边际代价是 \(\alpha_t\)，换回 \(\eta_c\)~kWh 储电、价值
\(\eta_c\nu_t\)，故仅当 \(\alpha_t\le\eta_c\nu_t\) 时充电；多放 1~kWh 可替代 \(\alpha_t\) 的购电成本，
代价是储电量减少 \(1/\eta_d\)~kWh，故仅当 \(\alpha_t\ge\nu_t/\eta_d\) 时放电。当 \(E_t>0\) 时由互补
松弛得 \(\alpha_t=p_t\)（实测偏差 \(\le10^{-6}\)）；当 \(E_t=0\) 时 \(\alpha_t\le p_t\)。由此得

\begin{proposition}[阈值律]\label{prop:threshold}
在储电量界不活跃（\(\underline S<S_t<\overline S\)）的时段，最优充放策略由两条阈值刻画：
\(\alpha_t\le\eta_c\nu_t\) 时充电，\(\alpha_t\ge\nu_t/\eta_d\) 时放电，而
\begin{equation}
  \eta_c\nu_t\ <\ \alpha_t\ <\ \nu_t/\eta_d
  \label{eq:q1-deadband}
\end{equation}
时 \(c_t=d_t=0\)，构成“死区”。
\end{proposition}

死区宽度随往返效率下降而变宽：两条阈值之比为 \((\nu_t/\eta_d)/(\eta_c\nu_t)=1/(\eta_c\eta_d)=1.2346\)。
即\textbf{放电时段的并网点电量边际价值须高于充电时段的 1.2346 倍，跨时段套利才有利可图}；在购电
时段该判据化为放电电价高于充电电价 1.2346 倍，此即式~\eqref{eq:ch1-arb} 的盈亏平衡比。

\textbf{适用边界。} 命题~\ref{prop:threshold} 在“纯套利子问题（变量仅 \(c_t,d_t,S_t\)）且储电量界
不活跃“时精确成立；界活跃时须回到式~\eqref{eq:q1-kkt} 并计入界乘子；功率界活跃时 \(z^c_t\) 或
\(z^d_t\) 不再为零，简单阈值关系失效。数值验证分两层：144 个时段上按 \(\alpha_t\) 判定充／放／死区的
\textbf{违反次数为 0}；界不活跃的 97 个时段子集（充电槽 42、放电槽 36）上违反次数同样为 0，
两条互补松弛恒等式最大残差为 0。实测 \(\nu_t\in[0.4711,1.1305]\)、\(\alpha_t\in[0.3713,1.2561]\)~元/kWh
均\textbf{时变}，故不存在“全天统一的电价阈值”；正确表述是两级结构——先由 LP 对偶确定时变水价值曲线，
再逐时段比较 \(\alpha_t\) 与两条阈值 \evid{C-Q1-02}。

\begin{figure}[htbp]
  \centering
  \includefig[width=0.60\linewidth]{fig03_q1_dual_price.pdf}
  \caption{最优策略的机制解释（三子图共享时间轴：光伏与负载功率、电价与储电量（右轴标出上下界
            1200/10800~kWh）、充放电量；对象：附件 1 典型日 144 槽）}
  \label{fig:dualprice}
\end{figure}

\subsection{无储能基线与成本分解}

无储能时 \(c_t=d_t=s_t=0\)，由式~\eqref{eq:q1-balance} 得闭式
\(E_t=\Delta t\max(L_t-PV_t,0)\)，据此定义基线

\begin{equation}
  C^{base}=\sum_t p_t\Delta t\max\bigl(L_t-PV_t,0\bigr),\qquad
  \sum_t p_tE_t=C^{base}+C^{batt},
  \label{eq:q1-decomp}
\end{equation}

其中 \(C^{batt}=\sum_t p_t\bigl(\Delta t\min(L_t-PV_t,0)+c_t-d_t+s_t\bigr)\) 为储能与弃光带来的成本
增量（典型日取负值即费用下降）。式~\eqref{eq:q1-decomp} 是\textbf{逐项恒等式}而非两个独立最优值
之和，\(C^{batt}\) 的最小化仍需求解式~\eqref{eq:q1-model}。基线有三个变体，\textbf{不得混用}
（表~\ref{tab:cbases}），否则全文多个“节省率”互不可比。

\begin{table}[htbp]
  \centering
  \caption{无储能基线 \(C^{base}\) 的三个变体及其适用问题（单位：元）}
  \label{tab:cbases}
  \footnotesize
  \begin{tabularx}{\linewidth}{C{1.0cm} L{3.6cm} L{3.0cm} C{1.6cm} Y}
    \toprule
    变体 & 口径 & 电价来源 & 适用问题 & 数值 \\
    \midrule
    V1 & 附件 1 典型日，单日周期 & 附件 1（固定价） & 问题 1 & 48\,052.0466 \\
    V2 & 334 天窗口，逐日独立 & 附件 1（固定价） & 问题 2、3 & 16\,407\,319.63 \\
    V3 & 334 天窗口，逐日独立 & 附件 4（波动价） & 问题 4 & 17\,026\,782.86 \\
    \bottomrule
  \end{tabularx}
\end{table}

\subsection{求解结果与机制}

在附件 1 典型日上求解式~\eqref{eq:q1-model}：全天购电量 \(59\,482.6990\)~kWh（\calref{C-02}）、
购电费 \(35\,126.9486\)~元，相对 \(C^{base}=48\,052.0466\)~元（变体 V1）节省
\(12\,925.0980\)~元，\textbf{节省率 \(26.90\%\)}（\calref{C-04}）\evid{C-Q1-03}；
充电量 \(20\,740.6661\)~kWh、放电量 \(16\,799.9396\)~kWh，二式互为倒数——
\(c/d=1.2346=1/(\eta_c\eta_d)\)，其倒数 \(d/c=0.8100=\eta_c\eta_d\) 即往返效率本身（\calref{C-07}）；
最优弃光 \(0.0000\)~kWh（\calref{C-06}）\evid{C-Q1-01}；储电量全程位于
\([1200.0000,10\,800.0000]\)~kWh，两端均为 \(6000.0000\)~kWh（严格日周期）。由
式~\eqref{eq:q1-decomp} 代入上述数值得 \(C^{batt}=-12\,925.0980\)~元，与目标函数直接相减一致。
表~\ref{tab:q1-buy}、表~\ref{tab:q1-soc} 按题面表 1、表 2 格式给出指定时段与 4~h 块的结果
（与 \texttt{result1.xlsx} 逐格一致）。

\begin{table}[htbp]
  \centering
  \caption{微网在指定时间段的购电量及全天的购电量和购电费（问题 1，单位：\unit{kWh} / \unit{元}）}
  \label{tab:q1-buy}
  \small
  \begin{tabular}{cccccc}
    \toprule
    时间段 & 购电量 & 时间段 & 购电量 & 时间段 & 购电量 \\
    \midrule
    \(10{:}00\)--\(10{:}10\) & 0.0000 & \(12{:}00\)--\(12{:}10\) & 480.4124 & \(14{:}00\)--\(14{:}10\) & 0.0000 \\
    \(16{:}00\)--\(16{:}10\) & 445.4317 & \(18{:}00\)--\(18{:}10\) & 531.8940 & \(20{:}00\)--\(20{:}10\) & 0.0000 \\
    \midrule
    \multicolumn{2}{c}{全天购电量（\unit{kWh}）：59\,482.6990} &
    \multicolumn{4}{c}{全天购电费（\unit{元}）：35\,126.9486} \\
    \bottomrule
  \end{tabular}
  \\[1ex]
  {\footnotesize 表注：时段按题面表述取区间左端时刻；\texttt{result1.xlsx} 沿用模板既有的时段末端
  标签，其数值与本文槽位按\textbf{位置}一一对应（槽 \(j\) 写入第 \(j+1\) 个数据行），标签仅作显示。}
\end{table}

\begin{table}[htbp]
  \centering
  \caption{储能设备在指定时间段的充放电量及 \(0{:}00\) 和 \(24{:}00\) 的储电量（问题 1，单位：\unit{kWh}）}
  \label{tab:q1-soc}
  \small
  \begin{tabular}{cccccc}
    \toprule
    时间段 & 充电量 & 放电量 & 时间段 & 充电量 & 放电量 \\
    \midrule
    \(0{:}00\)--\(4{:}00\)   & 4500.0000 & 0.0000    & \(4{:}00\)--\(8{:}00\)   & 833.3333  & 6365.8412 \\
    \(8{:}00\)--\(12{:}00\)  & 4787.9643 & 1702.9970 & \(12{:}00\)--\(16{:}00\) & 5286.0352 & 91.1014 \\
    \(16{:}00\)--\(20{:}00\) & 0.0000    & 5780.1319 & \(20{:}00\)--\(24{:}00\) & 5333.3333 & 2859.8681 \\
    \midrule
    \multicolumn{2}{c}{\(0{:}00\) 储电量：6000.0000} &
    \multicolumn{4}{c}{\(24{:}00\) 储电量：6000.0000} \\
    \bottomrule
  \end{tabular}
  \\[1ex]
  {\footnotesize 表注：每段含 36 个 10~min 槽；6 段之和与逐槽求和一致——充电 20\,740.6661~kWh、
  放电 16\,799.9396~kWh（\calref{C-07}）。}
\end{table}

\begin{figure}[htbp]
  \centering
  \includefig[width=0.54\linewidth]{fig10_strategy_mechanism.pdf}
  \caption{最优策略的机制解释（三子图共享时间轴：光伏与负载功率（\unit{kW}）、
            电价与储电量（\unit{元/kWh}、\unit{kWh}，储电量右轴标出上下界 1200/10800~\unit{kWh}）、
            充放电量（\unit{kWh}）；对象为附件 1 典型日的 144 个 10~\unit{min} 时段，
            竖直虚线连接光伏盈余时段、充电时刻与低价时段）}
  \label{fig:mechanism}
\end{figure}

图~\ref{fig:mechanism} 显示充电时段与光伏盈余时段、低价时段对齐，放电时段与高价时段对齐，
两组时段只\textbf{部分重合}，说明主机制是面向价格的跨时段套利，而非“把午间盈余搬到傍晚”。
可行性复核：供需平衡残差 \(1.49\times10^{-13}\)~kWh、储电量递推残差 \(9.09\times10^{-13}\)~kWh、
日周期偏差 \(\le10^{-12}\)~kWh；越界、同时充放与负购电的时段数均为 0；费用相对复算误差
\(\le10^{-6}\)（\S\ref{sec:audit}）。

\subsection{容量与功率的边际价值、效率口径}

在保持 10\%--90\% 放电深度（\(\underline S=0.1E_{cap}\)、\(\overline S=0.9E_{cap}\) 同步移动）下扫描
额定容量，并单独扫描功率上限（完整档位见 \appref{tab:scan-app}{E 容量与功率扫描明细}）：

\begin{equation}
  \begin{aligned}
    E_{cap}:12\,000\to28\,800\ \text{kWh} &\Rightarrow \Delta C=-6272.42\ \text{元}\ (-17.86\%),\\
    \overline P:5000\to10\,000\ \text{kW} &\Rightarrow \Delta C=-166.36\ \text{元}\ (-0.47\%),\\
    \text{仅解除 }\overline S &\Rightarrow -22.14\%\ (27\,351.55\ \text{元}),\qquad
    \text{仅解除 }\overline P &\Rightarrow -1.66\%\ (34\,544.43\ \text{元}).
  \end{aligned}
  \label{eq:q1-scan}
\end{equation}

界活跃统计：充电功率界在 \(19/144\) 个时段被打满，放电功率界 \(0/144\)；储电量上界活跃
\(21/144\)、下界活跃 \(8/144\)。故\textbf{容量摆幅是主要瓶颈，充放电功率仅在部分时段同时构成约束}，
两者量级相差一个数量级；与 \chapref{chap:verify} 的 Sobol 全局灵敏度一致（详见
\appref{tab:sobol-app}{E Sobol 灵敏度指数}）。

\textbf{效率口径。} 题面“充放电效率为 90\%”未区分充、放方向，故须显式声明口径
（主口径的选择理由见假设 A2）。以 \(\eta_c,\eta_d\) 的四种读法重解，得节省额

\begin{equation}
  \Delta C:\quad
  \underbrace{12\,925.10}_{0.9/0.9\ (\text{主口径})},\quad
  \underbrace{14\,250.99}_{\text{往返 }0.9},\quad
  \underbrace{14\,635.52}_{\text{仅充 }0.9},\quad
  \underbrace{13\,883.44}_{\text{仅放 }0.9}\quad(\text{元}),
  \label{eq:q1-eff}
\end{equation}

极差 \(1710.42\) 元，占主口径节省额的 \(13.2\%\)。故效率口径必须声明，否则同一模型可报出相差
八分之一的节省率。
